\section{Integral homotopy of the local sphere}
\label{sec:sphere}

We apply the coefficient theorem to the Morava $E$-Adams spectral
sequence and then recover integral homotopy from finite power cofibers.
The first step needs the horizontal vanishing theorem at a finite page,
rather than a bound on the cohomological dimension of the full
stabilizer.  The second step needs actual finiteness modulo $p$, rather
than a bound on torsion exponents.

\subsection{Descent after taking the cofiber of \texorpdfstring{$p$}{p}}

We recall the precise standard descent input.  For every prime $p$ and
positive height $n$, the commutative algebra $E_n$ is descendable in the
$K(n)$-local category.  The completed Amitsur complex for a spectrum $Y$
has terms
\[
 L_{K(n)}\bigl(Y\wedge E_n^{\wedge(q+1)}\bigr),\qquad q\geq0,
\]
and its totalization is $L_{K(n)}Y$.  Its Adams spectral sequence is
strongly convergent, and there are finite integers $R\geq2$ and $N\geq0$
such that its $R$th page vanishes in cohomological filtrations $s>N$.
For fixed $p,n$, these integers may be chosen independently of $Y$ and
of the internal degree.

These assertions are the $k=n$ case of
\citet[Theorem~4.6 and Proposition~4.13]{heard2023}; the identification
with the $K(n)$-local $E_n$-Adams spectral sequence is made in
\citet[Remark~4.14]{heard2023}.  The descendability statement follows
from the Hopkins--Ravenel smash product theorem and its preservation
under localization, as in \citet[Theorem~4.18 and
Proposition~10.10]{mathew2016}; the finite-page vanishing and convergence
are \citet[Corollary~4.4]{mathew2016}.  The fixed-point construction and
continuous-cohomology comparison are
\citet[Theorem~1\textup{(iii)--(iv)} and Appendix~A]{devinatzHopkins2004}.
None of these statements imposes a lower bound on $p$ relative to $n$.

Taking $Y=S/p$, the resulting spectral sequence is
\begin{equation}\label{eq:sphere-modp-descent}
 E_2^{s,t}
   =H^s_{\cts}\bigl(\mathbb G_n,\pi_t(E_n/p)\bigr)
 \ \Longrightarrow\
 \pi_{t-s}L_{K(n)}(S/p).
\end{equation}
Here $S/p$ is the finite Moore spectrum.  Thus the usual continuous
Amitsur comparison is compatible with smashing by $S/p$; equivalently,
apply the finite-dual form of
\citet[Theorem~1\textup{(iv)}]{devinatzHopkins2004} to $D(S/p)$.
This gives the indicated coefficient module, since $E_n$ is even and
$p$ is injective on its homotopy.  No commutative multiplication on the
Moore spectrum is required.

\begin{proposition}\label[proposition]{prop:sphere-modp}
For every prime $p$, every $n\geq1$, and every $j\in\mathbb Z$, the group
\[
 \pi_jL_{K(n)}(S/p)
\]
is finite.
\end{proposition}

\begin{proof}
Every entry on the second page of \eqref{eq:sphere-modp-descent} is finite
by \Cref{cor:coeff-morava}.  The entries on all later pages are
subquotients of those entries.  Choose $R,N$ as in the descent theorem.
Only $0\leq s\leq N$ can contribute to the limiting page, and the entries
contributing to the stem $j$ have $t=j+s$.  There are thus only finitely
many finite limiting entries in that stem.  Strong convergence supplies
a separated exhaustive filtration of the abutment. Above filtration
$N$, all successive quotients vanish, so that tail is constant;
separatedness makes the constant tail zero. The remaining filtration
is finite, with finite successive quotients, so the abutment is finite.  The same argument covers negative
$j$, and no evaluation or splitting of additive extensions is needed.
\end{proof}

\subsection{Finite power cofibers and integral recovery}

The mod-$p$ criterion for finite type appears for invertible
$K(2)$-local spectra in \citet[Proposition~3.1]{liFiniteType2021}.
We give the derived-complete form needed here, including the compact
Nakayama step. The argument applies in all integer degrees and imposes
no connectivity hypothesis.

For any spectrum $X$ and integer $a\geq1$, write
$X/p^a=\operatorname{cofib}(p^a:X\to X)$, with the quotient transition
maps $X/p^{a+1}\to X/p^a$.  We use derived $p$-completeness in the form
\[
 X\simeq\operatorname*{holim}_a X/p^a.
\]
This formulation does not assume ordinary completeness of its individual
homotopy groups.

\begin{lemma}\label[lemma]{lem:sphere-completion}
Let $X$ be a derived $p$-complete spectrum.  If $\pi_j(X/p)$ is finite
for every $j\in\mathbb Z$, then $\pi_jX$ is finitely generated over
$\Zp$ for every $j\in\mathbb Z$.  Conversely, finite generation of all
$\pi_jX$ over $\Zp$ implies finiteness of all $\pi_j(X/p)$.
\end{lemma}

\begin{proof}
The cofiber exact sequence gives, for every $j$,
\begin{equation}\label{eq:sphere-cofiber-short}
 0\longrightarrow \pi_jX/p\pi_jX
 \longrightarrow \pi_j(X/p)
 \longrightarrow (\pi_{j-1}X)[p]
 \longrightarrow0.
\end{equation}
Both end terms are killed by $p$.  Thus every finite middle term is a
finite $p$-group; we do not require it to be killed by $p$ itself.
Factoring $p^{a+1}$ as $p\circ p^a$ gives exact triangles
\[
 X/p^a\longrightarrow X/p^{a+1}\longrightarrow X/p.
\]
Their homotopy exact sequences inductively show that
$\pi_j(X/p^a)$ is a finite $p$-group for every $a\geq1$ and every $j$.

For fixed $j$, the inverse system of these finite groups is
Mittag--Leffler: its decreasing images in any one finite target
eventually stabilize.  Its first derived inverse limit therefore
vanishes.  Applying the Milnor exact sequence, in the adjacent degree
as well, and derived completeness gives
\begin{equation}\label{eq:sphere-finite-limit}
 \pi_jX\ \cong\ \varprojlim_a\pi_j(X/p^a).
\end{equation}
For clarity, vanishing of the derived limit requires no surjectivity of
the transition maps.  The map $1-\mathrm{shift}$ on the product of the
finite groups is surjective: every finite set of its defining equations
can be solved backwards, and compactness of the product supplies a
simultaneous solution.  This is also a direct proof of the required
lim-one vanishing.

Put $M=\pi_jX$.  The presentation \eqref{eq:sphere-finite-limit} makes
$M$ a compact Hausdorff pro-$p$ group with its canonical continuous
$\Zp$-module structure.  On each finite target $\pi_j(X/p^a)$, this scalar structure is
defined by reducing a $p$-adic integer modulo a power of $p$ that kills
that target; the definitions are compatible with the transition
maps.  Multiplication by $p$ on $M$ is continuous, so $pM$ is compact
and closed.  By \eqref{eq:sphere-cofiber-short}, the quotient $M/pM$ is
finite.

Choose $a_1,\ldots,a_d\in M$ lifting generators of $M/pM$ over
$\mathbb F_p$.  For $z\in M$, successively express
\[
 z_\ell=\sum_{i=1}^d b_{i,\ell}a_i+pz_{\ell+1},
 \qquad z_0=z,\qquad b_{i,\ell}\in\{0,\ldots,p-1\}.
\]
The scalars $c_i=\sum_{\ell\geq0}p^\ell b_{i,\ell}$ converge in $\Zp$.
The residual term after $L$ steps is $p^Lz_L$.  It tends to zero in
$M$: any fixed finite target in \eqref{eq:sphere-finite-limit} is
killed by a fixed power of $p$, which kills the image of $p^Lz_L$ for
all sufficiently large $L$, independently of $z_L$.  Therefore
$z=\sum_i c_i a_i$.  The continuous homomorphism
\[
 \Zp^d\longrightarrow M,\qquad (c_i)\longmapsto\sum_i c_i a_i
\]
is surjective.  This proves finite generation as an abstract
$\Zp$-module, and not merely density of a finitely generated subgroup.

Conversely, if every $\pi_jX$ is finitely generated over $\Zp$, then
both end terms of \eqref{eq:sphere-cofiber-short} are finite.  For the
$p$-torsion term, use that a submodule of a finite module over the
Noetherian ring $\Zp$ is finite as a module, and here it is killed by
$p$.  The middle term is consequently finite as well.
\end{proof}

\begin{proof}[Proof of \Cref{thm:main}]
Let $X=L_{K(n)}S_p^\wedge$.  A $K(n)$-local spectrum is derived
$p$-complete when $n>0$.  To see this directly, the fiber of
$X\to\operatorname*{holim}_aX/p^a$ is
\[
 \operatorname*{holim}(\cdots\xrightarrow{p}X\xrightarrow{p}X)
 \simeq F(S[1/p],X).
\]
The spectrum $S[1/p]$ is $K(n)$-acyclic, so this mapping spectrum
vanishes by locality.  This proves the required completeness without a
connectivity assumption.

The map $S\to S_p^\wedge$ is a mod-$p$ equivalence, hence a
$K(n)$-equivalence: its cofiber has invertible multiplication by $p$,
whereas multiplication by $p$ is zero on Morava $K$-theory.  Exactness
of localization therefore identifies
\[
 X\simeq L_{K(n)}S,\qquad X/p\simeq L_{K(n)}(S/p).
\]
The latter spectrum has finite homotopy groups in every degree by
\Cref{prop:sphere-modp}.  Applying \Cref{lem:sphere-completion} proves the sphere assertion.

For the finite-spectrum formulation, consider the class of finite
$p$-local spectra $F$ for which every
$\pi_j L_{K(n)}(F\wedge S/p)$ is finite. It contains the sphere by
\Cref{prop:sphere-modp} and is closed under suspensions, finite cofiber
sequences, and retracts, by exactness and the homotopy long exact
sequence. It therefore contains every finite $p$-local spectrum.
The spectrum $L_{K(n)}F$ is derived $p$-complete by the same locality
argument, and its cofiber of $p$ is $L_{K(n)}(F\wedge S/p)$.
Applying \Cref{lem:sphere-completion} objectwise proves the full
finite-spectrum conclusion of \Cref{thm:main}.
\end{proof}

\subsection{Ranks and the height-one calculation}

Finite generation gives an abstract decomposition
\[
 \pi_jL_{K(n)}S_p^\wedge
 \cong\Zp^{\,r_{p,n,j}}
   \oplus\bigoplus_{e\geq1}
          (\mathbb Z/p^e)^{m_{p,n,j,e}},
\]
with finitely many nonzero torsion multiplicities in each degree.  The
argument above establishes the existence of these finite invariants;
it does not determine the torsion multiplicities in general.

\begin{corollary}\label[corollary]{cor:sphere-ranks}
For every prime $p$, every $n\geq1$, and every $j\in\mathbb Z$, the free
rank is
\begin{equation}\label{eq:sphere-ranks}
 r_{p,n,j}
   =\#\left\{I\subseteq\{1,\ldots,n\}:
            \sum_{i\in I}(1-2i)=j\right\}.
\end{equation}
Equivalently, these ranks are the graded dimensions of the exterior
$\Qp$-algebra on generators in degrees $-1,-3,\ldots,-(2n-1)$.
\end{corollary}

\begin{proof}
By \citet[Theorem~A]{bssw2024}, the rationalized homotopy ring is the
indicated exterior algebra.  For a finitely generated $\Zp$-module its
$\Qp$-dimension after inverting $p$ is its free rank.  The exterior
monomials are indexed by subsets $I$ and have the degrees in
\eqref{eq:sphere-ranks}.  This proves the formula.
\end{proof}

\begin{proposition}\label[proposition]{prop:sphere-height-one}
Let $p$ be odd.  Then, for every $j\in\mathbb Z$,
\[
 \pi_jL_{K(1)}S_p^\wedge\cong
 \begin{cases}
  \Zp, & j=0\text{ or }j=-1,\\
  \mathbb Z/p^{\,1+\nu_p(t)},
       & j=2t-1,\quad0\neq t\in(p-1)\mathbb Z,\\
  0, & \text{otherwise}.
 \end{cases}
\]
The displayed groups include all additive extensions.
\end{proposition}

\begin{proof}
At height one, Morava $E$-theory is $p$-completed periodic complex
$K$-theory, with
\[
 \pi_*KU_p^\wedge=\Zp[\beta^{\pm1}],\qquad |\beta|=2.
\]
The stabilizer is
$\Zp^\times=\mu_{p-1}\times(1+p\Zp)$, and the Adams operation
$\psi^a$ acts on $\beta^t$ by $a^t$.
Taking homotopy fixed points by the finite group $\mu_{p-1}$ is exact
on these coefficient modules, because $p-1$ is a unit in $\Zp$.
The resulting Adams summand, denoted $B$, has
\[
 \pi_{2t}B=
 \begin{cases}\Zp,&(p-1)\mid t,\\0,&(p-1)\nmid t,
 \end{cases}
 \qquad \pi_{2t+1}B=0.
\]
The element $1+p$ topologically generates $1+p\Zp$.  Its two-term
continuous resolution, or the standard Adams-operation fiber sequence
\citep[\S\S1--2]{hopkinsKOne}, gives
\[
 L_{K(1)}S_p^\wedge\longrightarrow B
 \xrightarrow{\,\psi^{1+p}-1\,}B.
\]
On a retained coefficient in degree $2t$, the map is multiplication
by $(1+p)^t-1$.  For $t=0$ it is zero, giving $\Zp$ in degrees $0$
and $-1$.  For $t\neq0$ it is injective, and its cokernel is cyclic of
order determined by
\begin{equation}\label{eq:sphere-height-one-valuation}
 \nu_p\bigl((1+p)^t-1\bigr)=1+\nu_p(t).
\end{equation}
Indeed, if $p\nmid a$, binomial expansion gives
$(1+p)^a-1\equiv ap\pmod{p^2}$.  If $v_p(z)\geq1$ and $p$ is odd,
the term $pz$ in $(1+z)^p-1$ has valuation $1+v_p(z)$ and every other
term has strictly greater valuation.  Iterating this observation proves
\eqref{eq:sphere-height-one-valuation} for positive $t$.  For negative
$t$, the identity
\[
 (1+p)^{-t}-1=-\frac{(1+p)^t-1}{(1+p)^t}
\]
shows that the valuation is unchanged by replacing $t$ with $-t$.

The cokernel for a nonzero retained $t$ lies in degree $2t-1$; there is
no kernel in degree $2t$.  Coefficients in other degrees vanish.  The
fiber exact sequence has at most one contributing kernel or cokernel
in each stem, so the stated groups involve no unresolved additive
extension.
\end{proof}
